\documentclass[10pt,a4paper]{amsart}
\usepackage[margin=19mm]{geometry}
\usepackage{amsmath,amssymb,mathtools}
\usepackage[scr=rsfs,cal=euler]{mathalfa}
\usepackage{xfrac}
\usepackage[unicode,hidelinks,bookmarks=false]{hyperref}
\newcommand{\ie}{i.e.\ }
\newcommand{\cf}{cf.\ }
\newcommand{\resp}{respectively\ }
\newcommand{\Subsection}{\subsection}
\providecommand{\nobreakdash}{\nobreak\hbox{-}}
\setlength{\emergencystretch}{2em}
\multlinegap0pt
\usepackage{bm}
\usepackage{bbm}
\usepackage{dsfont}
\usepackage{xspace}
\usepackage{cancel}
\usepackage{mathtools}
\usepackage{amsthm}
\makeatletter
\@ifundefined{theorem}{\newtheorem{theorem}{Theorem}}{}
\makeatother
\def\P{\mathcal{P}}
\title[Strongly regular configurations]{Strongly regular configurations}
\author{Mari\'en Abreu, Martin Funk, Vedran Kr\v{c}adinac, Domenico Labbate}
\date{Source: arXiv:2104.04880v2 (CC BY 4.0)}
\begin{document}
\maketitle

\noindent\emph{Source and licence.} Strongly regular configurations. Version arXiv:2104.04880v2 (CC BY 4.0), distributed under the Creative Commons Attribution 4.0 International licence, as recorded in the arXiv licence metadata for this exact version and shown on the abstract page https://arxiv.org/abs/2104.04880v2. Full rights notice: licenses/arxiv-2104.04880-CC-BY-4.0.txt.\par\smallskip

\noindent\textit{Editorial scope.} Counted unit: the Theorem whose source label is \texttt{triangle} in the article (source0.tex lines 488--495, source label \texttt{triangle}), delivered together with the article's own complete original proof of it (source0.tex lines 497--527, one \texttt{proof} environment of 31 lines). No mathematics is written, repaired or re-derived: statement and proof are the article's own text line for line. Required context retained uncounted: none. Every definition, display and label of those lines is retained unchanged. Omitted from this package and not redistributed: 1--487, 496--496, 528--1240. Nothing inside a retained excerpt is omitted or altered: every retained body is delivered exactly as the article's lines are written, so no omission receipt of any kind accompanies this package. Citations: the retained excerpts carry no citation command at all, so no bibliography entry of the article is transcribed and no reference is redistributed. Reference adaptation: none was needed and none was made; every reference inside the delivered unit resolves inside it, so no unresolved reference or citation is produced. Printed slips kept literal: no printed word, symbol, hypothesis or number of the counted statement or of its proof was altered. Layout and class: the article's own class is \texttt{amsart} with packages including amssymb, amsfonts, amsmath, amsthm, color, hyperref, graphicx, amsaddr; the wrapper is the lane's pinned \texttt{amsart} editorial wrapper at 10pt on A4, which restates the theorem-like environments the retained text needs and adds the article's own macro definitions that the retained text needs (\texttt{\textbackslash P}), with extra wrapper packages (bm, bbm, dsfont, xspace, cancel, mathtools). The delivered numbering is the unit's own. Assets: no retained span contains a figure environment, a graphics-inclusion command, a PDF-page inclusion, a table or a TikZ picture; no image file is copied into this package, the package declares no asset dependency and no image file is redistributed with it. Licence: version arXiv:2104.04880v2 of this article is distributed under the Creative Commons Attribution 4.0 International licence, as recorded in the arXiv licence metadata for this exact version and shown on the abstract page https://arxiv.org/abs/2104.04880v2. A full rights notice is delivered at licenses/arxiv-2104.04880-CC-BY-4.0.txt. Characters: every retained excerpt is pure ASCII, so the delivered engine, XeLaTeX with its default Latin Modern text font, reports no missing character.
\begin{theorem}\label{triangle}
Let $\P$ be a projective plane of order $n\ge 5$ and $A$, $B$, $C$
be three non-collinear points. By deleting all points on the lines
$AB$, $AC$, $BC$ and all lines through the points $A$, $B$, $C$, there
remains a strongly regular $(v_k;\lambda,\mu)$ configuration with $v=(n-1)^2$, $k=n-2$,
$\lambda=(n-4)^2+1$, and $\mu=(n-3)(n-4)$. This configuration is not an
$(\alpha,\beta)$-geometry.
\end{theorem}

\begin{proof}
The number of points and lines in the remaining configuration is
$v=n^2+n+1-3-3(n-1)=(n-1)^2$ and they are of degree $k=n-2$. Let
$P$ and $Q$ be two remaining points that are collinear, i.e.\
are not on a line of $\P$ through $A$, $B$ or $C$. Then the points
non-collinear with $P$ are the remaining points on the
lines $AP$, $BP$, $CP$. There are $3(n-2)$ such points, and as many
for $Q$. The points non-collinear with both $P$ and $Q$ are
the intersections of one of the lines $AP$, $BP$, $CP$ with one of
the lines $AQ$, $BQ$, $CQ$; there are $6$ such points. By
inclusion-exclusion, the number of points in the remaining
configuration collinear with both $P$ and $Q$ is $\lambda=(n-1)^2-2
-6(n-2)+6=(n-4)^2+1$. If the points~$P$ and~$Q$
are non-collinear, a similar count shows that the number of
points in the remaining configuration collinear with both $P$ and $Q$ is
$\mu=(n-3)(n-4)$.

Let $(P,\ell)$ be a non-incident point-line pair of the remaining
configuration. We now count the points on $\ell$ collinear with $P$.
Let $A'$, $B'$, $C'$ be the intersections of $BC$, $AC$, $AB$ with~$\ell$.
These are the deleted points of~$\ell$. If the lines $AA'$, $BB'$, $CC'$
are concurrent, $P$ lies on $0$, $1$ or $3$ of these lines. Then there are
$n-5$, $n-4$ or $n-2$ points on~$\ell$ collinear with~$P$. In this case, the points
$A$, $B$, $C$, $A'$, $B'$, $C'$ and the common point of $AA'$, $BB'$, $CC'$
form a Fano subplane, so this can only occur if $n$ is even or $\P$ is non-Desarguesian.
On the other hand, if the lines $AA'$, $BB'$, $CC'$
are not concurrent, $P$ lies on $0$, $1$ or $2$ of these lines and there are
$n-5$, $n-4$ or $n-3$ points on~$\ell$ collinear with~$P$. In both cases there are three
possibilities for the number of points on~$\ell$ collinear with $P$, so the configuration
is not an $(\alpha,\beta)$-geometry.
\end{proof}

\end{document}
